
\medskip

\begin{enumerate}
\item \textbf{Réponse a.}

En effet, la fonction $f$ est dérivable sur $\R_{+}^{*}$, en tant que produit et somme de fonctions dérivables sur cet intervalle. Pour tout $x$ réel strictement positif, on a :

$f'(x)= 1 \times \ln(x) + x \times \dfrac{1}{x} - 1 + 0 = \ln(x) + \dfrac{x}{x} - 1 = \ln(x) +1-1 = \ln (x)$
\item \textbf{Réponse c.}
On a $f(x) = x^2 - x^2\ln (x)$.

$\bullet~~$On a $\displaystyle\lim_{x \to 0} x^2 = 0$ ;

$\bullet~~$On sait (croissances comparées) que  $\displaystyle\lim_{x \to 0} x^2 \ln (x) = 0$ ;

Donc par somme de limites $\displaystyle\lim_{x \to 0}f(x) = 0$.
%On a : $\lim\limits_{x \to +\infty} \ln(x) = +\infty$, donc, par limite de la somme, $\lim\limits_{x \to +\infty} \left[1 - \ln(x)\right] = -\infty$.
%
%Comme par ailleurs on a : $\lim\limits_{x \to +\infty} x^2 = +\infty$, par limite du produit, on en conclut :
%
%$\lim\limits_{x \to +\infty} g(x) = -\infty$.
\item \textbf{Réponse d.}

Quel que soit $x \in \R,\: f(x) = x\left(x^2 - 0,9x - 0,1\right)$.

On a donc $f(x) = 0 \iff \left\{\begin{array}{l c l}
x&=&0\\
x^2 - 0,9x - 0,1&=&0
\end{array}\right.$

$0$ est donc une solution et pour l'équation du second degré le nombre 1 est une racine évidente ; le produit des racines étant égal à $-0,1$ l'autre racine est donc $- 0,1$.

L'équation $f(x) = 0$  a donc trois solutions : $-0,1 \:;\ 0$ et 1.
%$f$ sur $\R$. $f$ est dérivable sur $\R$, en tant que fonction polynôme. Pour tout $x$ réel, on a :

%$f'(x) = 3x^2 - 0,9 \times 2x - 0,1 = 3x^2 - 1,8x - 0,1$.
%
%Cette dérivée est un polynôme du second degré, dont le discriminant est :
%
%$\Delta = (-,18)^2 -4\times (3) \times (-0,1) = 4,44 >0$.
%
%Il y a donc deux racines réelles distinctes : $x_1 = \dfrac{1,8 + \sqrt{4,44}}{6} $ et $x_2 =\dfrac{1,8 - \sqrt{4,44}}{6}$.
%
%			Comme le coefficient dominant est positif, cela signifie que $f'(x)$ est positif, sauf pour $x$ dans l'intervalle $[x_1~;~x_2]$ (car $x_1 < x_2$).
%
%			$f$ est donc strictement croissante sur les intervalles $ ]-\infty~;~x_1] $ et $ [x_2~;~+\infty[$ et strictement décroissante sur $ [x_1~;~x_2] $.
%
%			On calcule des valeurs approchées des images de $x_1$ et $x_2$ (à la calculatrice) :
%
%			$f(x_1) \approx -0,27$ et $f(x_2) \approx -0,44$.
%
%			En combinant ces images et les variations, on peut dire que, $f(x_1)$ est le maximum de $f$ sur l'intervalle $ ]-\infty~;~x_2] $, et comme ce maximum est strictement négatif, l'équation n'a pas de solution sur cet intervalle.
%
%			Par ailleurs, $\lim\limits_{x \to +\infty} f(x) = \lim\limits_{x \to +\infty} x^3\left(1 + - \dfrac{0,9}{x} - \dfrac{0,1}{x^2}\right) = +\infty$.
%
%			Sur l'intervalle $[x_2~;~+\infty[$, $f$ est une fonction continue (c'est un polynôme), strictement croissante et 0 est une valeur intermédiaire entre $f(x_2) <0$ et $\lim\limits_{x \to +\infty} f(x) = +\infty$. D'après le théorème des valeurs intermédiaires pour les fonctions strictement croissantes, l'équation $f(x) = 0$ admet donc une unique solution sur cet intervalle.
%
%			Globalement, l'équation admet donc exactement une unique solution sur $\R$.
			\item \textbf{Réponse c.}

			On dérive la proposition, en pensant bien à utiliser la formule pour dériver une fonction composée ($x\longmapsto H(2x)$ a pour dérivée $x \longmapsto 2H'(2x)$) et on vérifie.

			Si $K_\textbf{c} : x \longmapsto \dfrac{1}{2} H(2x)$, alors,

			$\forall x \in \R, K_\textbf{c}' (x) = \dfrac{1}{2}\times 2H'(2x) = H'(2x) = h(2x)$, car $H$ est une primitive de $h$ sur $\R$

			et donc on a : $\forall x \in \R, K_\textbf{c}' (x) = k(x)$.

			On en déduit donc que $K_\textbf{c}$ est une primitive de $k$ sur $\R$.

			\item \textbf{Réponse b.}

			Dérivons la fonction $f$, qui est dérivable sur $\R$, en tant que produit de fonctions dérivables sur cet ensemble.

			$\forall x \in \R, \quad 1 \times \mathrm{e}^x + x \times \mathrm{e}^x = (1 + x)\mathrm{e}^x$.

			On a donc $f'(1) = (1 + 1)\mathrm{e}^1 = 2\mathrm{e}$.

			Par ailleurs $f(1) = 1\times \mathrm{e}^1 = \mathrm{e}$.

			L'équation réduite de la tangente à la courbe de $f$ au point d'abscisse 1 est donc :
			$\aligned[t] y &= f'(1) (x - 1) + f(1)\\ &= 2\mathrm{e} (x - 1) + \mathrm{e} \\&= 2\mathrm{e}x - 2\mathrm{e} + \mathrm{e} \\y&= 2\mathrm{e}x - \mathrm{e}\endaligned$

			\item \textbf{Réponse d.}

			Résolvons cette inéquation :

			$ \aligned [t] (0,2)^n < 0,001 &\iff \ln\left(0,2^n\right) < \ln(0,001) \text{\qquad car ln est strictement croissante sur }\R^{*+}\\
			&\iff n\ln(0,02) < \ln(0,001)\\
			&\iff n > \dfrac{\ln(0,001)}{\ln(0,2)} \text{\qquad car } \ln(0,2)< 0
			\endaligned$

			On a $\dfrac{\ln(0,001)}{\ln(0,2)} \approx 4,3$, donc les solutions à cette inéquation sont les entiers naturels supérieurs ou égaux à 5.
		\end{enumerate}

\bigskip

